\subsection*{Limitations}
\label{sec:limitations}

The experimental evidence is bounded in several ways.
(i)~All ten datasets satisfy $n \leq 600$ and $d \leq 60$; behaviour on
high-dimensional or large-scale data is not tested.
(ii)~Only two importance methods (PAI and NRS) are compared head-to-head;
the remaining three (ReliefF, mutual information, variance) serve as
baselines but are not the focus of Pattern~B/C analysis.
(iii)~Only the Vietoris--Rips filtration with bottleneck distance is used;
alternative filtrations (\v{C}ech, alpha) and diagram distances
(Wasserstein-$p$, landscape) may behave differently.
(iv)~The systematic search did not include Scopus or Web of Science
(unavailable via programmatic API in this research environment); relevant
papers indexed only in those databases may be missed.
(v)~Citation counts (retrieved September 2025) are a snapshot and
will drift.
(vi)~NRS uses a fixed $k=7$ and PAI uses a single filtration configuration
($d_{\max}=1$, threshold $=2.0$) throughout; no hyperparameter ablation was
performed on these settings (Section~\ref{sec:ablation}).

\section{Conclusion}
\label{sec:conclusion}

We have surveyed two decades of work bringing modern mathematical and
machine-learning frameworks into granular computing's core pipeline, organised
around a single thesis: such imports tend to exhibit one or more of three structural failure
patterns---reduction to a known granular quantity (Pattern~A), dominance of
label-orthogonal geometry (Pattern~B), or self-optimising evaluation (Pattern~C).

Our survey makes the following contributions:
\begin{enumerate}
  \item \textbf{Taxonomy.}  A systematic mapping
        of the topological research into two non-communicating branches
        (abstract topology and computational TDA), with an analysis of the gap
        between them and the conditions under which bridging is possible without
        triggering Pattern~A.
  \item \textbf{Saturation assessment.}  A taxonomy of granule
        geometries (Table~\ref{tab:granule_tax}) showing that six deformation
        axes of the neighbourhood ball have each been explored by multiple
        groups, with diminishing returns on new directions.
  \item \textbf{Negative-result experiments.}  A dataset of $4{,}975$
        attribute-importance measurements and $30{,}100$ downstream
        classification records, released per-fold, documenting where a
        prominent candidate method (persistence-based importance) fails and why.
        Theorem~\ref{thm:pai_bound} bounds the effect and
        Section~\ref{sec:patB} tests it empirically.  A further ablation
        (Table~\ref{tab:knn_scale}) shows that the single metric favouring
        the method reverses its verdict when that metric's own scale
        parameter is varied, and reports that our proposed explanation for
        the effect failed its own test.
  \item \textbf{Methodology critique.}  A data-driven critique of the
        3-NN/SVM-UCI evaluation protocol, supported by a full-text coding of
        reporting practice in 22 papers, 13 of them a majority of the
        Axis-1 corpus (Table~\ref{tab:practice}), released with a supporting
        quotation for every cell, and by evidence
        that on most standard benchmarks the margin between the best reduct
        and the all-features baseline is not statistically distinguishable
        from zero
        (seven of ten datasets gain $\leq 2$ pp; a Holm-corrected paired
        Wilcoxon finds the gain significant on only three), rendering
        between-method comparisons on those benchmarks uninformative.
  \item \textbf{Trend analysis with failure-pattern diagnostics.}
        An assessment of three active research directions---incremental
        feature selection, three-way decision integration, and
        neural-augmented granular computing---identifying the precise
        conditions under which each avoids Patterns~A--C and constitutes a
        non-trivial contribution (Section~\ref{sec:trends}).
  \item \textbf{Diagnostic checklist and reporting standard.}
        Actionable guidance for researchers: three checks before implementing
        a new GrC extension (Section~\ref{sec:checklist}), and six minimum
        reporting requirements for experimental papers
        (Section~\ref{sec:standard}).
\end{enumerate}

The nine open problems in Section~\ref{sec:open} represent the frontier
where progress is possible---including the theoretically open questions of
reduct quality bounds for incremental algorithms and data-driven cost matrices
for three-way decision---requiring new theoretical machinery (cross-axis
synthesis, high-dimensional acceleration, multi-scale separation axioms)
rather than geometric variations on existing methods.

\section*{Declaration of competing interest}
The authors declare no known competing financial interests or personal
relationships that could have appeared to influence the work reported in this
paper.  References~\cite{DaiTT2024IFT} and~\cite{Nguyen2026MHMG} are the
authors' own prior work; the former is
identified as such in Section~\ref{sec:abstract_topo} and is assessed against
the same diagnostic checklist applied to every other proposal surveyed here.

\section*{Acknowledgements}
No funding to declare.  All experiments were conducted on a personal laptop.

\section*{Data and Code Availability}
All experimental data and code are deposited at
\url{https://doi.org/10.5281/zenodo.XXXXXXX} [DOI to be finalised before
submission].  The deposit includes:
\begin{itemize}[leftmargin=*,itemsep=1pt]
  \item \emph{Raw measurements.}
        \texttt{raw\_importance.csv} ($4{,}975$ rows),
        \texttt{raw\_downstream.csv} ($30{,}100$ rows),
        \texttt{raw\_structure.csv} ($995$ rows), and
        \texttt{raw\_knn\_scale.csv} ($5{,}970$ rows, the metric-scale
        ablation of Section~\ref{sec:patC}).
  \item \emph{Experiment scripts.} \texttt{importance.py},
        \texttt{run\_e1.py}, \texttt{datasets.py}, \texttt{run\_structure.py},
        \texttt{run\_knn\_scale.py}, \texttt{analyze\_knn\_scale.py}.
  \item \emph{Presentation and audit.} \texttt{gen\_tables.py}, which
        regenerates every data table from the raw files, together with
        \texttt{verify\_tables.py}, \texttt{check\_dois.py} and
        \texttt{check\_authors.py}, which re-derive every printed number from
        the raw data and re-resolve every reference against Crossref.  A
        reviewer can run them to confirm that no number in this paper was
        transcribed by hand and that no citation is unresolvable.
  \item \emph{Review artefacts.} \texttt{literature\_log.csv} (the 52-paper
        systematic-review log) and \texttt{fulltext\_coding.csv} (the
        per-paper reporting-practice coding of Section~\ref{sec:significance},
        with the supporting quotation for every judgement).
\end{itemize}
